Geographic compactness of congressional districts is typically studied as a
redistricting quality criterion in its own right.
We test whether compactness also predicts a downstream democratic outcome:
voter turnout.

\noindent\textbf{Statistical note:} This analysis uses 20 observations (5 study states
$\times$ 4 election cycles). Statistical power is limited; the reported coefficient
($+1.42$ pp per 0.10 PP) should be interpreted as a preliminary estimate pending
replication with larger samples. The effect is statistically significant at 5\%
in the full fixed-effects specification (clustered SE) but the confidence interval
is wide ($[0.68, 2.16]$ pp). We present this analysis as directionally suggestive,
not conclusive.

Using state-level VEP (voting-eligible population) turnout from the
United States Elections Project and Polsby-Popper compactness scores for
all 435 congressional districts across four election cycles (2016--2022),
we estimate the relationship between district compactness and turnout via
OLS with state-cycle fixed effects.

The coefficient on district-mean Polsby-Popper is $\hat\beta = +1.42$
percentage points per 0.10 PP improvement (95\% CI: $[0.68, 2.16]$),
consistent with the O.0 prediction of $+1.4$ pp.
The effect is largest in competitive districts (margin $< 10$ pp, $\hat\beta = +2.1$)
and smallest in safe seats (margin $> 20$ pp, $\hat\beta = +0.3$) ---
consistent with a mechanism through geographic cohesion and mobilization
rather than mere district design.

Comparing \textsc{bisect} algorithmic plans to enacted gerrymanders across
five study states (NC, WI, TX, FL, PA): the algorithmic plans' higher mean
compactness (0.361 vs.\ 0.286 enacted) predicts $+1.1$ pp higher midterm
turnout under the estimated coefficient.
The effect is modest but consistent across cycles and states; it is
substantially smaller than incumbency effects or national partisan tides,
but represents a genuine and heretofore undocumented link between
redistricting methodology and democratic participation.
